% =====================================================================
%  DANE PRACY — jedyne miejsce, w którym wpisuje się metadane.
%  Tytuł, streszczenie i słowa kluczowe MUSZĄ być identyczne z USOS-APD
%  (Regulamin studiów §35 ust. 5). Rok = rok złożenia pracy (APD: 18.12.2026).
% =====================================================================
\newcommand{\pracaAutor}{Krystian Przepiórka}
\newcommand{\pracaOpiekun}{dr inż. Paweł Urbański}
\newcommand{\pracaKierunek}{Inteligentna Elektronika}
\newcommand{\pracaSpecjalnosc}{}   % pusta = wiersz „Specjalność” pominięty
\newcommand{\pracaRok}{2026}
\newcommand{\pracaTytul}{Opracowanie, wykonanie oraz badania miniaturowego
  analizatora EDS w~technologii MEMS do zastosowań kosmicznych}
% Tytuł EN z repozytorium W:\VAULT\Inzynierka\latex (body.tex, \titlen)
\newcommand{\pracaTytulEN}{Development, implementation and testing of
  a~miniature MEMS-based EDS analyzer for space applications}

% Streszczenie i abstract: wersje robocze z repozytorium
% (frontmatter/streszczenie.tex, abstract_en.tex) — do uzupełnienia o wyniki
% i ujednolicenia z tekstem w USOS-APD.
\newcommand{\pracaStreszczenie}{%
  Celem pracy jest opracowanie, wykonanie oraz badania miniaturowego
  analizatora EDS (\angielski{Energy Dispersive X-ray Spectroscopy})
  w~technologii MEMS (\angielski{Micro-Electro-Mechanical Systems}),
  przeznaczonego do zastosowań kosmicznych. W~ramach pracy zaprojektowano
  źródło elektronów wykorzystujące zjawisko emisji polowej z~nanorurek
  węglowych (CNT) wraz z~elektrodą ekstrakcyjną oraz elektrodą umożliwiającą
  przyspieszanie wiązki elektronów. Przeanalizowano możliwość zastosowania
  magnetycznego układu ogniskowania wiązki elektronów w~celu poprawy
  parametrów wiązki oraz zwiększenia efektywności wzbudzania
  charakterystycznego promieniowania rentgenowskiego. Wykonano prototyp
  układu, a~jego działanie zweryfikowano eksperymentalnie z~wykorzystaniem
  spektrometru rentgenowskiego. Oceniono możliwość wykorzystania technologii
  MEMS do budowy kompaktowych analizatorów materiałowych przeznaczonych do
  zastosowań kosmicznych.
  \TODO{dopisać 2--3 zdania z~wynikami liczbowymi; tekst 1:1 z~USOS-APD (200--300 słów)}}
\newcommand{\pracaSlowaKluczowe}{%
  MEMS, spektroskopia EDS, emisja polowa, nanorurki węglowe,
  ogniskowanie magnetyczne, zastosowania kosmiczne}
\newcommand{\pracaAbstract}{%
  The aim of this thesis is to develop, build and test a~miniature EDS
  (Energy Dispersive X-ray Spectroscopy) analyzer manufactured in MEMS
  (Micro-Electro-Mechanical Systems) technology, intended for space
  applications. The work covers the design of an electron source based on
  cold field emission from carbon nanotubes (CNT), combined with an
  extraction electrode and an accelerating electrode. The feasibility of
  using a~magnetic electron-beam focusing system is analyzed in order to
  improve beam parameters and increase the efficiency of characteristic
  X-ray excitation. A~working prototype is built and verified
  experimentally using an X-ray spectrometer. The feasibility of using
  MEMS technology to build compact material analyzers for space
  applications is evaluated.
  \TODO{add 2--3 sentences with key results; text identical with USOS-APD}}
\newcommand{\pracaKeywords}{%
  MEMS, EDS, field emission, carbon nanotubes, magnetic focusing,
  space applications}

% Metadane PDF
\hypersetup{pdfauthor={\pracaAutor}, pdftitle={\pracaTytul},
            pdfsubject={Praca dyplomowa inżynierska, PWr WEFiM}}
